\newpage
\section{面向初学者的 LuaLaTeX：Lua 脚本入门}

\subsection{LaTeX 中的 Lua 简介}

\tkzEngine{\LUALATEX} 是 LaTeX 的一个变体，它将 Lua 集成为脚本语言。这样可以：

\begin{itemize}
  \item 执行高级计算。
  \item 处理文本与数据。
  \item 在文档中生成动态内容。
\end{itemize}

\subsection{在 LaTeX 文档中使用 Lua}

在 \tkzNamePack{tkz-elements} 中，我只使用两种把代码集成到 \tkzname{Lua} 的方式，如下：

\begin{itemize}
  \item  使用 \tkzMacro{lualatex}{directlua} 的 Lua 代码块。注意\footnote{加载同名宏包后可使用 \tkzEnv{luacode}{luacode} 环境。也可以使用 \tkzEnv{tkz-elements}{tkzelements} 环境。}

  Lua 代码可以使用 \tkzMacro{lualatex}{directlua} 在文档中直接执行：

  |\directlua{tex.print(math.exp(1))}| -->  \directlua{tex.print(math.exp(1))}

  \item  加载外部 Lua 文件：

  你可以编写独立的 Lua 脚本，并用 \tkzFct{lua}{dofile}、\tkzFct{lua}{loadfile} 或 \tkzFct{lua}{require} 来加载。这三个函数各有优劣，视情况而定；本教程先重点介绍前两个。

 \emph{dofile：}   立即执行一个 Lua 文件。等价于 loadfile(filename)()

\emph{loadfile：}    将 Lua 文件载入内存（不立即执行）；返回一个待执行的函数。

\emph{require：} 将 Lua 模块加载一次（通过 package.loaded）并执行；使用 package.path 查找文件。


本文档中给出了示例，可参见：

 [\ref{sec:work_organization}] |\directlua{dofile("lua/sangaku.lua")}|


\end{itemize}

\tkzRHand 在 \tkzEngine{\LUALATEX} 下如何在文件中写注释。有几种情形需要注意：在 |directlua| 指令之外的 \LATEX{} 部分，始终使用 \% 符号。在 |\directlua| 指令内部，仍使用 \%。在外部文件 |fic.lua| 中，Lua 的单行注释符号是 |--|（两个连字符）。

在 \tkzEnv{luacode}{luacode} 环境中，需要使用两个连字符。

下一小节介绍 \TeX{} 与 LuaTeX{} 在特殊字符和类别码（catcode）上的一些差异。

\subsection{\tkzNamePack{Lua\TeX} 中的特殊字符与类别码} % (fold)
\label{sub:special_characters_and_catcodes}

\tkzHand{\textcolor{red}{警告}}：
使用 \tkzNamePack{Lua\TeX} 时，会暴露出 Lua 与 \TeX{} 在字符处理方式（尤其是特殊字符）以及类别码（catcode）的分配与解释上的深层差异。

\subsubsection{\LaTeX{} 中的类别码是什么？} % (fold)
\label{ssub:what_are_catcodes}

在 \LaTeX{} 中，每个字符都关联一个\emph{类别码}（catcode），它决定 \TeX{} 在输入处理时如何解释该字符。


\begin{tabular}{lll}
\toprule
 \textbf{字符} & \textbf{类别码} & \textbf{含义} \\
 \midrule
 \verb|\| & 0 & 转义字符（开始一个控制序列） \\
\verb|{| & 1 & 开始分组 \\
\verb|}| & 2 & 结束分组 \\
\verb|#| & 6 & 宏参数 \\
\verb|^| & 7 & 上标（数学模式中） \\
\verb|%| & 14 & 注释字符 \\
 \bottomrule
 \end{tabular}


这些类别码可以在宏或文档中动态改变。这种灵活性很强大，但在与 Lua 交互时也有风险。

% subsubsection what_are_catcodes（结束）

\subsubsection{Lua 不理解类别码} % (fold)
\label{sub:lua_does_not_understand_catcodes}

Lua 将字符串解释为字节序列或 UTF-8 码点序列。它\emph{没有类别码的概念}。当 Lua 通过 \tkzMacro{lualatex}{directlua} 向 \TeX{} 发送内容时，由 \TeX{} 根据当前类别码设置进行解释。

这种差异会导致常见问题：

\begin{itemize}
  \item 特殊字符（如 \verb|#|、\verb|%| 或 \verb|^|）未被转义，可能破坏 \LaTeX{} 解析器。
  \item 若 \verb|\| 使用不当或被重定义，控制序列可能被误解。
  \item 从 Lua 插入的分组定界符（\verb|{| 和 \verb|}|）可能被错误解析，尤其是通过 \tkzFct{lualatex}{tex.sprint} 注入宏时。
\end{itemize}

% subsubsection lua_does_not_understand_catcodes（结束）

\subsubsection{典型问题示例} % (fold)
\label{ssub:typical_problems}

\begin{itemize}
  \item \textbf{井号（\#）}：从 Lua 打印 \verb|\mycmd{#1}| 时，若 \verb|#| 的类别码不是 6 或未正确转义，可能报错。
  \item \textbf{换行符（\texttt{\textbackslash n}）}：Lua 用 \verb|\n| 表示换行，但它在 TeX 输入中无效，除非通过 \tkzMacro{latex}{scantokens} 传递。
  \item \textbf{百分号（\%）}：\% 符号对 (La)TeX 和 Lua 都是特殊符号，但方式截然不同。在 (La)TeX 中它是默认注释字符；在 Lua 中它是模式匹配中的“魔法”字符之一。（Lua 的其他“魔法”字符是 |^$()[]*+-?]|。）

\end{itemize}


\subsection{Lua 与 LaTeX 的交互}

\tkzEngine{\LUALATEX} 允许你与 LaTeX 交互，特别是通过：

\begin{itemize}
\item \code{|tex.print()|}：在 LaTeX 中显示文本。
\item \code{|tex.sprint()|}：插入内容而不添加额外空格。\\
\end{itemize}

\begin{mybox}
\begin{verbatim}
|\directlua{tex.sprint("\\textbf{Bold text generated by Lua}")}|
\end{verbatim}
\directlua{tex.sprint("\\textbf{Bold text generated by Lua}")}
\end{mybox}


\subsection{变量}

\begin{mybox}
  在 Lua 中，变量是存储某个值的命名标签。它的类型由所赋的值动态决定，其作用域（全局或局部）决定了它在程序中何处可用。
\end{mybox}

我们来分析刚才这句话。有四个要点很重要：“命名标签”、“值的类型”、“动态地”和“作用域”。


\subsubsection{命名标签}

你给变量取的名字必须遵循一定规则：

\begin{itemize}
  \item 可以包含字母（a–z、A–Z）、数字（0–9）和下划线（\_）。
  \item 不能以数字开头。
  \item 区分大小写。
  \item 某些 Lua 关键字（如 if、then、function、local 等）是保留字，不能用作变量名。
\end{itemize}

\subsubsection{数据类型}
\label{ssub:types}
 可以给变量赋哪些类型的值？下面列出最基本的几种：

\begin{itemize}
    \item \textbf{nil}：未定义的值。
    \item \textbf{boolean}：\texttt{true} 或 \texttt{false}。
    \item \textbf{number}：双精度浮点数（实数）。
    \item \textbf{string}：字符串。
    \item \textbf{table}：关联数据结构（键–值）和数组。
    \item \textbf{function}：匿名或命名函数。
\end{itemize}

\subsubsection{动态类型}

类型由所赋的值决定。因此 \code{|x = 5|} 时 \code{x} 的类型是 \code{number}，而 \code{|x = true|} 时它的类型变为 \code{boolean}。

\tkzRHand 前面曾警告过 \tkzNamePack{tkz-elements} 所保留的类与变量。例如，你不能使用 \code{circle}，因为 $circle = ...$ 会导致无法再使用 \tkzClass{circle} 这个类。同样，$C= 1$ 会导致表 \tkzVar{circle}{C} 的内容丢失。

\subsubsection{作用域与块}

\begin{enumerate}

  \item  \code{\textbf{块}}\\
在 Lua 中，块是被当作一个整体处理的语句序列。Lua 中的块由特定关键字界定。它们对于组织代码、定义局部变量的作用域以及控制执行流程至关重要。

下面是 Lua 中定义块的主要方式：

\begin{itemize}
  \item \code{|do ... end|}。显式地创建一个块。在 do...end 块中声明的局部变量只能在该块内访问。

  \item \code{|if ... then ... else ... end|} 或 \code{|if...then...elseif...then...end|}。then 与 end（或 else 与 end，或 elseif 之后的 then）之间的语句构成一个按条件执行的块。

\item \code{|for...do...end|} 或 \code{|while...do...end|}。循环体（do 与 end 之间）是一个被重复执行的块。

\item \code{|function...end|}。函数体（function 与 end 之间）是一个包含函数代码的块。
\end{itemize}


\item \code{\textbf{作用域}}\\
Lua 主要有两种作用域：全局或局部。变量的作用域决定了它在代码中何处可访问。默认情况下，一个被赋值而未显式声明（用 local）的变量是全局的。建议优先在函数或块内只声明局部变量，做法是在变量名前加 \code{local}。

在块内声明为 \code{local} 的变量在块外不可访问。

示例：

\begin{mybox}
\begin{verbatim}
\directlua{
 do
  local x = 10
  tex.print(x) % Accessible here
end
 tex.print(x) % Error! x is out of scope here}
\end{verbatim}
\end{mybox}
\end{enumerate}

\subsection{值}

我们来看看关于主要值类型需要了解的内容。

%   \begin{enumerate}[label=(\greek*)]

\subsubsection{\code{nil}}


  \begin{itemize}

    \item  \code{nil} 是最基本的，并且自身就是一种数据类型。\code{nil} 用于表示变量尚未赋值，或某个值被显式清除。

\begin{mybox}
  \begin{verbatim}
    \directlua{
    tex.print(tostring(no_assigned))}
    and
    tex.print(type(no_assigned)) gives
  \end{verbatim}
\directlua{tex.print(tostring(no_assigned))}
\directlua{tex.print(type(no_assigned))}
\end{mybox}

\item 表与 \code{nil}。\\

给某个键赋 \code{nil} 会将该条目从表中删除。可以用 \code{nil} 删除一个表。若该表没有其他引用，这样做可以释放内存。

  \end{itemize}

\subsubsection{布尔值}

\begin{itemize}
  \item 在 Lua 中只有两个布尔值：\code{true} 和 \code{false}。
  \item 布尔值属于 boolean 数据类型。
  \item 使用逻辑运算符（\code{and}、\code{or}、\code{not}）对它们进行操作，以构成更复杂的布尔表达式。
  \item 与许多其他语言不同，在 Lua 中只有 false 和 nil 在布尔上下文中（例如 if 条件中）被视为假。若想区分 nil 与 false，需要更显式地判断：
\begin{mybox}
\begin{verbatim}
  if v == false then
   ...
  elseif v == nil then
   ...
  else
   ...
  end
\end{verbatim}
\end{mybox}
\end{itemize}

\subsubsection{数值}


\begin{itemize}
  \item 单一数值类型。Lua 只有一种数值类型来表示实数（双精度浮点数）。这意味着所有数字，无论看起来像整数还是带小数点，都被同等对待。

\item Lua 使用标准的 64 位浮点表示。
\item 尽管只有一种数值类型，Lua 仍能高效处理整数值。

\item  算术运算符。Lua 支持标准算术运算符。值得注意的是 |%|（取模）——除法的余数，以及 |//|（整除）——相除后将结果向下取整到最近的整数。你可以使用 math 库中定义的函数。

\item 可以用 \code{tostring} 函数把数字转换为字符串。例如：a = 10、b = "1"、c = tostring(a) .. b。最终 |tex.print(c)| 得到 $101$。\\ 注意：在某些情况下，\code{加法}或\code{拼接}等运算会改变其中一个操作数的类型以得到结果，但也有许多特殊情形。

\end{itemize}

\subsubsection{字符串}

\begin{enumerate}[label=(\alph*)]
  \item 构造

 你可以在 Lua 中用单引号（'...'）、双引号（"..."）或双方括号（[[...]]）定义字符串字面量，但这几种方式有差异。如果字符串不含任何特殊字符，就不会有大问题。


第一个小问题是想在字符串中包含单引号或双引号。最简单的方法是用双方括号定义字符串。你可以在双引号定义的字符串中包含单引号，反之亦然。

\begin{tkzexample}[latex=.4\textwidth]
\directlua{tex.print('a string with " ')}
\end{tkzexample}

如果你想稍微挑战一下，也可以在双引号定义的字符串中包含双引号。这一次，为了用 \TeX{} 中常见的做法，我们尝试写 \"，但这会报错。只需用 \string\" 使 \" 失效即可得到正确结果。

\begin{tkzexample}[latex=.4\textwidth]
\directlua{tex.print("a string with \string\" ")}
\end{tkzexample}



\item 拼接：\code{|..|}。一个有趣的例子是可以创建多个变量：
\begin{mybox}
  \begin{verbatim}
    local t= {}
    local i = 5
    t["x_" .. tostring(i)] = 11
    tex.print(t.x_5)
  \end{verbatim}

\end{mybox}


\item Lua 提供了强大的字符串库（string），其中有许多用于操作字符串的有用函数。下面是一些常用的：

\begin{itemize}

\item \code{string.char(n)},
\item \code{string.len(s)},

\item \code{string.sub(s, i [, j])},

\item \code{string.find(s, pattern [, init [, plain]])},

\item \code{string.gsub(s, pattern, repl [, n])},

\item \code{string.format(formatstring, ...)} 等。
\end{itemize}

\item \code{String} 库中的特殊或魔法字符\\

\begin{center}
  \begin{tabular}{@{}ll@{}}
  \toprule
  \textbf{字符} & \textbf{含义} \\
  \midrule
  \texttt{.} & 匹配任意字符 \\
  \texttt{\%} & 转义字符（用于魔法字符） \\
  \texttt{+} & 一次或多次重复 \\
  \texttt{-} & 零次或多次重复（最小匹配） \\
  \texttt{*} & 零次或多次重复（贪婪匹配） \\
  \texttt{?} & 零次或一次出现 \\
  \texttt{[ ]} & 字符类 \\
  \texttt{\^{}} & 字符串开头 \\
  \texttt{\$} & 字符串结尾 \\
  \bottomrule
  \end{tabular}
\end{center}

\item 字符类（位于 \texttt{\%} 之后）

\begin{center}
  \begin{tabular}{@{}ll@{}}
  \toprule
  \textbf{代码} & \textbf{含义} \\
  \midrule
  \texttt{a} & 字母（A–Z、a–z） \\
  \texttt{c} & 控制字符 \\
  \texttt{d} & 数字（0–9） \\
  \texttt{g} & 可打印字符（空格除外） \\
  \texttt{l} & 小写字母 \\
  \texttt{p} & 标点字符 \\
  \texttt{s} & 空白字符 \\
  \texttt{u} & 大写字母 \\
  \texttt{w} & 字母数字字符 \\
  \texttt{x} & 十六进制数字（0–9、A–F、a–f） \\
  \texttt{z} & 空字符（ASCII 0） \\
  \bottomrule
  \end{tabular}
\end{center}

\item 说明
\begin{itemize}
    \item 要匹配字面魔法字符（如 \texttt{.} 或 \texttt{*}），在其前面加 \texttt{\%}。
    \item 大写代码（\texttt{\%A}、\texttt{\%D} 等）匹配\emph{除}对应字符类之外的任意字符。
\end{itemize}

\item LuaLaTeX 中的 String.format —— 速查表
我选择把 \tkzname{Lua} 代码作为 |directlua| 指令的参数，这意味着你要留意任何可能想使用的特殊符号（\#、\%、\_ 等）。

\begin{center}
  \begin{tabular}{@{}ll@{}}
  \toprule
  \textbf{Lua 格式}   & \textbf{说明} \\ \midrule
  \texttt{\%d}      & 十进制整数   (7 或 7.0)                           \\
  \texttt{\%i}      & 整数或十六进制数的占位符         \\
  \texttt{\%f}      & 浮点数（固定精度）      \\
  \texttt{\%g}      & 浮点数（自动精度）  \\
  \texttt{\%s}      & 字符串占位符                            \\
  \texttt{\%e}      & 科学计数法（指数）               \\
  \texttt{\%\%}     & 百分号（转义）                      \\
  \texttt{\%-5d}    & 将整数左对齐，字段宽度为 5  \\
  \texttt{\%10.2f}  & 浮点数右对齐，宽度 10，保留 2 位小数 \\
  \texttt{\%.2f}    & 保留 2 位小数的浮点数  \\
  \texttt{\%g}      & 通用格式的浮点数    \\ \bottomrule
  \end{tabular}
\end{center}


* 基本格式
\begin{itemize}
  \item \texttt{string.format(\%d, 5)} 输出：\texttt{5}
  \item \texttt{string.format(\%f, 2.5)} 输出：\texttt{2.500000}
  \item \texttt{string.format(\%g, 1234567.89)} 输出：\texttt{1.23457e+06}
  \item \texttt{string.format(\%s, "z.A")} 输出：\texttt{z.A}
\end{itemize}

* LaTeX 中的数学表达式
你可以轻松地格式化幂、方程等数学表达式：
\begin{itemize}
  \item \texttt{string.format(\$ \%d \textasciicircum 3 = \%d \$, x, x\textasciicircum3)} 将输出类似 \texttt{\$8\textasciicircum 3 = 512\$} 的内容
\end{itemize}

* 特殊字符
\begin{itemize}
  \item \texttt{string.format(\%\%)} 输出：\texttt{\%}
\end{itemize}


\item 若 s = "example"，则 |tex.print(s:len())| 得到 $7$。
\end{enumerate}

\subsubsection{表}

\begin{itemize}

\item 定义\\
在 Lua 中，table 类型实现了关联数组，并且是唯一可用的数据结构。这意味着在两个实体之间建立了关联：第一个称为键（key），第二个称为值（value）。每个键关联唯一的值。键可以是整数、实数（少见但可能）、字符串、布尔值、表、函数，但不能是 nil。

  一个很常见的特殊情形是键为从 1 开始的连续整数的表，这种表简称为数组，此时键就是下标。


  \item 表的创建。\\
  创建表最简单的方法是用构造表达式 \{\}，但下面的例子使用了语法糖：

  \begin{mybox}
    \code{t = \{one = 1, two = 4, three = 9, four = 16, five = 25\}} \\
    键–值对之间用逗号或分号分隔。键在 = 号左边，值在右边。
   | t.one = 1; t.two = 4; etc.|
  \end{mybox}

  实际上，这等价于

  \begin{mybox}
    \code{t = \{["one"] = 1, ["two"] = 4, ["three"] = 9, ["four"] = 16, ["five"] = 25\}} \\
  \end{mybox}

  这是表最一般的定义形式。

   | t.["one"] = 1; t.["two"] = 4; etc.|

  使用语法糖对用户来说很方便，但它隐藏了一些要点。首先，字符串键的\code{点记法}只有在字符串键符合 Lua 标识符规则（字母数字和下划线，不以数字开头）时才有效。

  \item 表操作。\\
Lua 内置了用于表操作的函数，列在下表中。

\begin{table}[htbp]
\centering
\caption{Lua 中内置的表操作函数}
\begin{tabular}{@{}ll@{}}
\toprule
\textbf{函数} & \textbf{说明} \\
\midrule
\texttt{table.insert(t, [pos,] value)} & 在表 \texttt{t} 的位置 \texttt{pos} 插入一个值（默认：末尾）。 \\
\texttt{table.remove(t, [pos])}        & 从表 \texttt{t} 的位置 \texttt{pos} 删除元素。 \\
\texttt{table.sort(t [, comp])}        & 原地排序表 \texttt{t} 的元素。可选比较函数。 \\
\texttt{table.concat(t [, sep [, i [, j]]])} & 用分隔符 \texttt{sep} 拼接 \texttt{t} 中从 \texttt{i} 到 \texttt{j} 的元素。 \\
\texttt{table.unpack(t [, i [, j]])}   & 返回 \texttt{t} 中下标从 \texttt{i} 到 \texttt{j} 的元素。 \\
\texttt{table.pack(...)}          & 返回一个新表，把所有参数存为其元素。 \\

\bottomrule
\end{tabular}
\end{table}






\item 访问元素：可以用方括号 [] 加相应键来访问表元素。若键是合法的字符串标识符，也可以使用点记法（.）。
\item 删除元素：要从表中删除一个元素，只需把它的键赋值为 nil。
\end{itemize}

\subsubsection{函数}

在 Lua 中，函数具有一等地位，这意味着它们可以被赋值给变量、作为参数传给其他函数，以及作为函数的返回值。

 \begin{enumerate}[label=(\alph*)]
  \item 基本语法

\begin{tkzexample}[latex=.4\textwidth]
  \directlua{
    function cubic(x)
       return x^3
    end
    tex.print(cubic(5))}
\end{tkzexample}

\item 参数与返回值\\

函数可以接受任意数量的参数并返回多个值。下面的例子说明这种情形，并展示与 Lua 特殊字符相关的一些困难。

这涉及整数除法。最简单的做法是在外部文件 “lua/divide.lua” 中创建这个函数：

\begin{mybox}
\begin{verbatim}
 function divide_and_remainder(dividend, divisor)
   if divisor == 0 then
     tex.error("Error: Cannot divide by zero")
     return nil
   else
     local quotient = math.floor(dividend / divisor)
     local remainder = dividend % divisor
     return quotient, remainder
   end
 end
\end{verbatim}
\end{mybox}

\begin{verbatim}
\directlua{
 dofile("lua/divide.lua")
local pc = string.char(37)
local num1 = 17
local num2 = 5
local q, r = divide_and_remainder(num1, num2)
tex.sprint(string.format("The quotient of ".. pc.. "d divided by"
.. pc.." d is:".. pc.." d\\par", num1, num2, q))
tex.sprint(string.format("The remainder is: "..pc.."d\\par", r))
}
\end{verbatim}

\directlua{
 dofile("lua/divide.lua")
local pc = string.char(37)
local num1 = 17
local num2 = 5
local q, r = divide_and_remainder(num1, num2)
tex.sprint(string.format("The quotient of ".. pc.. "d divided by".. pc.."
 d is:".. pc.." d\\par", num1, num2, q))
tex.sprint(string.format("The remainder is: "..pc.."d\\par", r))
}

另一种做法是使用局部变量，例如：

\verb|local format = string.char(37) .. ".2f"|

然后 \verb| tex.print(string.format(format , ...))|

我们来看看函数定义在 \tkzMacro{lualatex}{directlua} 宏中的版本。

问题在于使用运算符 \%（取模运算符，整数除法后的余数）。在使用 string.format 函数以及 \%d 或 \%i（整数占位符）时也会遇到同样的问题。

这两种情形的一个解决办法是定义一个宏 \tkzname{pc} 来代替 \%。

\begin{verbatim}
 \makeatletter
 \let\pc\@percentchar
 \makeatother
\end{verbatim}

代码变为：


\begin{verbatim}
\directlua{
 function divide_and_remainder(dividend, divisor)
   if divisor == 0 then
     tex.error("Error: Cannot divide by zero")
     return nil
   else
     local quotient = math.floor(dividend / divisor)
     local remainder = dividend \pc divisor
     return quotient, remainder
   end
 end

local num1 = 17
local num2 = 5
local q, r = divide_and_remainder(num1, num2)

tex.sprint(string.format([[The quotient of \pc i divided by
 \pc i is: \pc i\\]], num1, num2, q))
tex.sprint(string.format([[The remainder is: \pc i]], r))}
\end{verbatim}

\item 变参函数\\

用 ... 来处理可变数量的参数：

第一个例子：

这是避开 \# 相关问题最简单的方法。

\begin{itemize}


  \item  \begin{tkzexample}[latex=.2\textwidth]
    \directlua{
    local function sum_all(...)
      local total = 0
      for _, value in ipairs{...} do
        total = total + value
      end
      return total
    end

    tex.print(sum_all(1, 2, 3, 4))
    }
  \end{tkzexample}


\item  第二个例子：

这里我使用 \tkzFct{utils}{table\_getn} 包中的一个新函数，它取代了旧的 Lua 函数 \tkzFct{lua}{table.getn}（用于获取表的大小）。

  \begin{tkzexample}[latex=.2\textwidth]
    \directlua{
    local function sum_all(...)
      local arg = {...}
      local total = 0
      for i = 1, utils.table_getn(arg) do
        total = total + arg[i]
      end
      return total
    end
    local result = sum_all(10, 20, 30, 40)
    tex.print(result)}
  \end{tkzexample}


\item  第三个例子：

只要稍下功夫（技巧），仍然可以使用 \#。

\begin{tkzexample}[latex=.2\textwidth]
 \bgroup
 \catcode`\#=12
 \directlua{
 local function sum_all(...)
   local arg = {...}
   local total = 0
   for i = 1, #arg do
     total = total + arg[i]
   end
   return total
 end
 local result = sum_all(10, 20, 30, 40)
 tex.print(result)}
 \egroup
\end{tkzexample}
\end{itemize}


\item 递归\\

Lua 支持递归函数：

\begin{tkzexample}[latex=.4\textwidth]
  \directlua{dofile("lua/fact.lua")}
  \newcommand*{\luafact}[1]{\directlua{
    tex.write(fact(\the\numexpr(#1)\relax))}%
  }
  $8! = \luafact{8}$
\end{tkzexample}



\item 方法（对象式语法）\\

在为表（对象）定义函数时，可以使用冒号语法（:），自动把对象本身作为第一个参数传入，通常称为 self。

下面是一个例子：


\begin{tkzexample}[latex = .2\textwidth]
  \directlua{
   z.A = point(0, 0)
   z.B = point(5, 0)
   L.AB = line(z.A, z.B)
   T.ABC = L.AB:half()
   _,_,z.C = T.ABC:get()
   tex.print(tostring(z.C))}
\end{tkzexample}

若不使用冒号语法，你需要显式地把对象作为第一个参数传入：

\begin{verbatim}
  T.ABC = L.AB.half(L.AB)
  _,_,z.C = T.ABC.get(T.ABC)
\end{verbatim}

\code{half} 和 \code{get} 是为对象表定义的方法。用冒号语法调用时，它们会自动接收所操作的实例（self）。
\end{enumerate}

\subsection{控制结构}


控制结构让你的程序能够做判断并重复执行动作。Lua 有一小套易用的控制结构，帮助你写出灵活而强大的代码。

\textbf{1. 条件语句}

当你希望代码根据条件执行不同动作时，使用 if、elseif 和 else。这样程序就能在不同选项之间“选择”。

\textbf{2. 循环}

循环用于重复执行动作：

当条件为真时，while 重复执行。
repeat ... until 一直重复直到条件为真（它至少执行一次）。
for 用于计数或遍历表中的元素。

\textbf{3. 局部变量与块}

当你用 local 创建变量时，它只存在于你写它的那个块内。块可以是函数、循环或 if 语句。

\textbf{4. 跳出与返回}

\emph{break} 让你提前结束循环。
return 从函数中返回一个值，或完全结束函数。
Lua 保持简洁，一旦学会这些基本结构，你就能在程序中编写各种逻辑！

   \begin{enumerate}
  \item \code{if then else}
  \begin{tkzexample}[latex = .35\textwidth]
    \directlua{
    x = math.sin(math.pi)
     tex.print("x = ",x)
       tex.print([[\\]])
    if math.abs(x) < 1e-12 then
    tex.print("$x = 0$")
    elseif x >0 then
     tex.print("$x > 0$")
     else
      tex.print("$x < 0$")
    end
    }
  \end{tkzexample}


\item \code{while}
\begin{tkzexample}[latex = .35\textwidth]
\directlua{
 i = 1
 while i <= 5 do
   tex.print(i)
   i = i + 1
 end}
\end{tkzexample}

\item \code{repeat}
\begin{tkzexample}[latex = .35\textwidth]
\directlua{
 x = 0
 repeat
   x = x + 2
   tex.print(x)
 until x == 6}
\end{tkzexample}

  \item 数值型 \code{for}：\\
 | for init,max/min value, increment do statement(s)  end|



\begin{tkzexample}[latex = .35\textwidth]
\directlua{
  numbers = {10, 20, 30, 40, 50}
  local nb = utils.table_getn(numbers)
  for i = 1, nb do
    local sep = (i < nb) and ", " or "."
    tex.sprint(numbers[i] .. sep)
  end}
\end{tkzexample}


  \item 泛型 \code{for}：\\
   |for i,v in ipairs(t) do tex.print(v) end|

  \begin{tkzexample}[latex = .35\textwidth]
 \directlua{
    t = {a = 1, b = 2}
    for k, v in pairs(t) do
      tex.print(k, v)
    end}
  \end{tkzexample}

  本包的一个示例：\\

\begin{tkzexample}[latex = .35\textwidth]
  \directlua{
  init_elements()
  z.A = point(0, 1)
  z.B = point(2, -1)
  z.C = point(1, 1)
  for k, v in pairs(z) do
    tex.sprint(k,": ","(".. tostring(v)..")")
    tex.print('\\\\')
  end}
\end{tkzexample}

  \item 关于 \code{for} 的小结
  \begin{center}
    \bgroup
    \small
  \captionof{table}{Lua 中 \texttt{ipairs} 与 \texttt{pairs} 的比较}\label{for:summary}

  \begin{tabular}{@{} lcc @{}}
  \toprule
   \textbf{特性} & \texttt{ipairs(t)} & \texttt{pairs(t)} \\
   \midrule
   遍历      & 数值键 \texttt{1..n}      & 所有键（字符串、数字等） \\
   顺序              & 有保证（1、2、3、\ldots）     & 无保证 \\
   停止于           & 序列中第一个 \texttt{nil}  & 所有键之后 \\
   适用场景           & 数组 / 列表                  & 表 / 字典 \\
   跳过非数值键  & \cmark                          & \xmark \\
   包含空洞     & \xmark                          & \cmark \\
   \bottomrule
   \end{tabular}
  \egroup
  \end{center}



\item \code{break} 与 \code{return}
  \begin{itemize}
  \item     break：退出循环。

\begin{tkzexample}[latex = .35\textwidth]
\directlua{  for i = 1, 10 do
    if i == 5 then break end
    tex.print(i)
  end}
\end{tkzexample}

\item    return：退出函数（或程序块），可带返回值也可不带。

下面是一个带有如下约束的例子：

\begin{enumerate}[label=(\roman*)]

\item   函数遍历一个数字列表。

\item   若找到偶数（value \% 2 == 0），立即返回它。

\item   若遇到 $-1$，用 break 结束循环，不立即返回值。

\item   若在遇到 $-1$ 之前未找到偶数，函数最后返回 nil。
\end{enumerate}


\begin{tkzexample}[latex=.25\textwidth]

   \makeatletter
   \let\pc\@percentchar
   \makeatother

\directlua{%
  function find_first_even(t)
    for i = 1, utils.table_getn(t) do
      local value = t[i]
      if value \pc 2 == 0 then
        return value
      elseif value == -1 then
        break
      end
    end
    return nil
  end

  local numbers = {1, 3, 5, -1, 4, 6}
  tex.print(tostring(find_first_even(numbers))) }
\end{tkzexample}

\end{itemize}

\end{enumerate}

\subsection{Lua 语法糖}
\label{sub:lua_sugar_syntax}
Lua 中的“语法糖”指那些让语言更简洁、更易读而不改变其根本含义的语法元素。下面是 Lua 中主要的语法糖形式：

   \begin{enumerate}
  \item  用 “:” 进行方法调用\\
 \code{ obj:method(arg)} 是 \code{ obj.method(obj, arg)} 的语法糖。

 “:” 隐式地把 obj 作为第一个参数（self）传入，这在面向对象编程中非常有用。


  \item  用 “.” 访问表字段\\
  \code{t.key} 是 \code{  t["key"]} 的语法糖。

  当键是合法标识符（不含空格、标点等）时，这样更易读。

  \item 简写表构造式\\
  \code{local t = \{1, 2, 3\}} 等价于：\code{ local t = \{[1]=1, [2]=2, [3]=3\}}

  并且：

  \code{local t = \{a = 1, b = 2\}} 是 \code{local t = \{["a"] = 1, ["b"] = 2\}} 的语法糖。

  \item 循环语法糖（例如 for k, v in pairs(t) do）\\
  尽管 pairs() 只是一个函数，但 Lua 用 for 遍历表的方式非常简洁易读。k 表示键，v 表示值。

\begin{mybox}
\begin{verbatim}
  for k, v in pairs(t) do
     ...
  end
\end{verbatim}
\end{mybox}


 \item 用逻辑表达式实现默认值的技巧\\
 x = a or b
 通常表示：“若 a 为真则用它，否则用 b”。非常适合设置默认值。

 \item do ... end 块\\
 这是一种创建局部作用域块的方式——用于限制变量的可见性。
\end{enumerate}

\subsection{示例：计算 1 到 10 的和}
\label{sub:calculating_the_sum}

Lua 可用于 LaTeX 本身无法高效处理的高级计算。

\vspace{1em}
\begin{tkzexample}[latex=.5\textwidth]
\directlua{
  local somme = 0
  for i = 1, 10 do
    somme = somme + i
  end
  tex.print("Sum of 1 to 10: " .. somme)}
\end{tkzexample}
%

\subsubsection{示例：斐波那契}

 \vspace{1em}
 计算第 $10$ 项

\vspace{1em}
\begin{tkzexample}[latex = .4\textwidth]
\directlua{
 function fibonacci(n)
   if n == 0 then
     return 0
   elseif n == 1 then
     return 1
   else
     return fibonacci(n - 1) + fibonacci(n - 2)
   end
 end}
    斐波那契数列中第 $10$ 项为
    \(u_{10}=\directlua{tex.print(
    fibonacci(10))}\)
\end{tkzexample}
\endinput
